%! AP Statistics — Unit 4, Lesson 4.5 — Group Activity
\documentclass[10pt]{article}
\usepackage{apstats-article}
\usepackage{apstats-boxes}

\begin{document}

\pageheader{Unit 4, Lesson 4.5}{Group Activity: Conditional Probability from Tables}
\namepartnerperiod

\vspace{0.1in}

\begin{scenariobox}[Context]{sky}
\small
A fitness researcher surveyed 180 athletes about their preferred sport and
how frequently they train. The results are shown below.

\vspace{6pt}
\renewcommand{\arraystretch}{1.4}
\begin{tabularx}{\linewidth}{@{} l X X X r @{}}
 & \textbf{Daily} & \textbf{3--4$\times$/week} & \textbf{1--2$\times$/week} & \textbf{Total} \\
\hline
\textbf{Running}  & 30 & 20 & 10 & 60 \\
\textbf{Swimming} & 24 & 18 & 18 & 60 \\
\textbf{Cycling}  & 18 & 22 & 20 & 60 \\
\hline
\textbf{Total}    & 72 & 60 & 48 & 180 \\
\end{tabularx}

\vspace{4pt}
One athlete is selected at random.
\end{scenariobox}

\vspace{0.1in}

% ── TIER R ────────────────────────────────────────────────────────────────────
\begin{tcolorbox}[colback=greenbg, colframe=greenacc, boxrule=0.9pt, arc=2mm,
    left=3mm, right=3mm, top=2mm, bottom=2mm,
    title={\bfseries\small Tier R --- Remediate: Identifying Numerators and Denominators},
    fonttitle=\small]
\small
For each conditional probability expression, identify the numerator (joint count)
and denominator (conditioning total) from the table. You do not need to compute yet.

\vspace{6pt}
\renewcommand{\arraystretch}{1.6}
\begin{tabularx}{\linewidth}{@{} l X X @{}}
\textbf{Expression} & \textbf{Numerator (which cell?)} & \textbf{Denominator (which total?)} \\
\hline
$P(\text{Running} | \text{Daily})$ & \blank{3cm} & \blank{3cm} \\
$P(\text{Daily} | \text{Swimming})$ & \blank{3cm} & \blank{3cm} \\
$P(\text{Cycling} | \text{1--2$\times$/week})$ & \blank{3cm} & \blank{3cm} \\
$P(\text{Swimming} | \text{3--4$\times$/week})$ & \blank{3cm} & \blank{3cm} \\
\end{tabularx}

\vspace{6pt}
Now compute each probability above. Show your division.\\[6pt]
$P(\text{Running}|\text{Daily}) = $ \blank{2cm} $\div$ \blank{2cm} $=$ \blank{3cm}\\[4pt]
$P(\text{Daily}|\text{Swimming}) = $ \blank{2cm} $\div$ \blank{2cm} $=$ \blank{3cm}
\end{tcolorbox}

\vspace{0.1in}

% ── TIER A ────────────────────────────────────────────────────────────────────
\begin{tcolorbox}[colback=sky, colframe=navy, boxrule=0.9pt, arc=2mm,
    left=3mm, right=3mm, top=2mm, bottom=2mm,
    title={\bfseries\small Tier A --- Approaching Proficiency: Calculate, Apply, and Interpret},
    fonttitle=\small]
\small
\begin{enumerate}[leftmargin=1.5em, itemsep=10pt]
  \item Compute the following conditional probabilities. Write the formula, substitute, and simplify.
        \begin{enumerate}[label=(\alph*), itemsep=6pt]
          \item $P(\text{Running} | \text{Daily})$: \blank{8cm}
          \item $P(\text{Daily} | \text{Running})$: \blank{8cm}
          \item Are these two values equal? What does that tell you? \writeline\writeline
        \end{enumerate}

  \item Apply the multiplication rule to verify your answer to (a).\\
        $P(\text{Daily}) \cdot P(\text{Running}|\text{Daily}) = $\blank{3cm}$\times$\blank{3cm}$ = $\blank{3cm}\\
        Does this equal $P(\text{Running} \cap \text{Daily})$ from the table? \blank{2cm}

  \item Compute $P(\text{Cycling} | \text{3--4$\times$/week})$ and interpret in a complete sentence.\\
        Calculation: \blank{6cm}\\
        Interpretation: \writeline\writeline
\end{enumerate}
\end{tcolorbox}

\vspace{0.1in}

% ── TIER E ────────────────────────────────────────────────────────────────────
\begin{tcolorbox}[colback=goldbg, colframe=goldacc, boxrule=0.9pt, arc=2mm,
    left=3mm, right=3mm, top=2mm, bottom=2mm,
    title={\bfseries\small Tier E --- Extension: Reverse Conditioning and Missing Values},
    fonttitle=\small]
\small
\begin{enumerate}[leftmargin=1.5em, itemsep=10pt]
  \item Compare $P(\text{Swimming} | \text{Daily})$ and $P(\text{Daily} | \text{Swimming})$.
        Compute both values, then explain in context why they are different and which
        question each answers.\\
        $P(\text{Swimming}|\text{Daily}) = $\blank{4cm} \quad $P(\text{Daily}|\text{Swimming}) = $\blank{4cm}\\
        Explanation: \writeline\writeline\writeline

  \item Suppose a new group of 30 athletes is added: 12 runners who train daily,
        10 swimmers who train 3--4$\times$/week, and 8 cyclists who train 1--2$\times$/week.
        Recompute $P(\text{Running} | \text{Daily})$ with the updated table.
        Does the value increase, decrease, or stay the same? Explain why.\\
        \writeline\writeline\writeline

  \item Create a conditional probability question that could be answered using this table,
        where the answer is greater than 0.5. Write the question, compute the answer,
        and explain why it makes sense in context.\\
        \writeline\writeline\writeline
\end{enumerate}
\end{tcolorbox}

\end{document}
